% 第6章：本章小结。
\FloatBarrier
\section{本章小结}
\label{sec:matching-assignment-summary}

本章开头区分了两种要求：只求配对最多，或要求每个对象都配上并使总成本最低。对第一种，二部图匹配与单位容量网络上的整数流一一对应，残量增广路径落到二部图上就是交替路径，于是匹配最大当且仅当不存在增广路径。这个判据比“不能再加边”强：极大匹配可能仍有增广路径，改进要允许中间顶点换伙伴。匈牙利算法只维护伙伴数组，逐个左侧顶点用深度优先搜索找增广路径，总时间$O(n+pm)$，其中$p$为较小一侧的顶点数；Hopcroft--Karp算法按层成批增广，得到$O(n+m\sqrt n)$。这些构造都依赖给定的两侧划分，不能直接用于含奇数长环的一般图。

对第二种，可行标号给出全部完整指派的共同下界，只用紧边的完美匹配达到这个下界，因而最优。KM算法在紧边上生长交替树，受阻时按最小松弛量调标，使下界升高、树继续生长；结束时匹配与标号构成精确证书，禁配也能直接处理，稠密方阵上为$O(r^3)$。Auction算法用价格协调争抢，维护$\varepsilon$互补松弛：最终指派的成本至多比最优值高$r\varepsilon$，只有在整数成本且$\varepsilon<1/r$时才保证精确最优；出价次数随成本范围与$\varepsilon$变化，总时间为$O(r^2+rB)$。

两类结果的核对方式也不同。最大匹配用单位容量网络中与配对数同值的割核对；KM算法的结果逐边核对标号可行与所选边取等号；Auction算法的结果只能核对最终价格下的$\varepsilon$条件，再据此给出误差界。矩形、禁配与允许不配对的情形，都要先补成方阵并把惩罚写进成本，求解后按原始成本重算。

标号和价格都是给顶点配上的数，用来为配对成本提供可核查的界：价格经$u_i=\min_j(c_{i,j}+p_j)$、$v_j=-p_j$就转成可行标号。下一章把这种做法推广为一般容量与供需网络上的节点势，最大流、匹配与指派也都能写成那里的最小费用流。
